\documentclass[12pt,twoside]{article}

\input{macros-sp26}
\newcommand{\thelecturenum}{2}

\title{CS336 Lecture 2}

\begin{document}

\handout{Lecture \thelecturenum: Resource Accounting (and PyTorch)}

\setlength{\parindent}{0pt}

\subsection*{Motivating the Accounting}

\textbf{Question:} What knowledge to take away from this lecture?

\begin{itemize}
    \item Mechanics: straightforward (PyTorch semantics).
    \item Mindset: resource accounting (remember to do it).
    \item Intuitions: get a sense of how resources are spent, no ML magic today.
\end{itemize}

\textbf{Question:} How long would it take to train a 70B parameter model on 15T tokens on 1024 H100s?

\begin{minted}{python}
    total_flops = 6 * 70e9 * 15e12
    h100_flop_per_sec = 1979e12 / 2
    mfu = 0.5
    flops_per_day = h100_flop_per_sec * mfu * 1024 * 60 * 60 * 24
    days = total_flops / flops_per_day
\end{minted}

\textbf{Question:} What's the largest model that can you can train on 8 H100s using AdamW?

\begin{minted}{python}
    h100_bytes = 80e9
    bytes_per_parameter = 12
    num_parameters = (h100_bytes * 8) / bytes_per_parameter
\end{minted}

\begin{remark}
    Activations are not accounted for (depends on batch size and sequence length), so this is an upper bound.
\end{remark}

This is a rough back-of-the-envelope calculation. But it gives you the flavor of napkin math one can quickly do to get a sense of resources.

\begin{minted}{python}
def get_memory_usage(x: torch.Tensor):
    return x.numel() * x.element_size()
\end{minted}

\subsection*{Tensor Basics}

Tensors are the basic object for storing everything (e.g. data, parameters, gradients, activations).

\begin{xample}
    Parameters of the DeepSeek v3.2 model.
\end{xample}

Each tensor has a rank, which is the number of dimensions.

\begin{minted}{python}
    x = torch.zeros(4)        # rank 1 tensor (vector) 
    x = torch.zeros(4, 8)     # rank 2 tensor (matrix) 
    x = torch.zeros(4, 8, 2)  # rank 3 tensor
\end{minted}

In Transformers, will see tensors of rank 4:

\begin{minted}{python}
    B = 32   # Batch size
    S = 16   # Sequence length
    H = 16   # Number of heads
    D = 64   # Hidden dimension per head
    x = torch.zeros(B, S, H, D)
\end{minted}

\subsection*{Tensor Memory}

Elements of tensors are generally floating point numbers. The fp32 data type (also known as single precision) is the default representation. Traditionally, in scientific computing, fp32 is the baseline; you could use double precision (fp64) in some cases. In deep learning, you can be a lot sloppier. 

\begin{minted}{python}
def get_promised_flop_per_sec(dtype: torch.dtype) -> float:
    """Return the peak FLOP/s for `device` operating on `dtype`."""
    if not torch.cuda.is_available():
        return 1
    properties = torch.cuda.get_device_properties(cuda_if_available())  
    if "A100" in properties.name:
        if dtype == torch.float32:
            return 19.5e12
        if dtype in (torch.bfloat16, torch.float16):
            return 312e12
        raise ValueError(f"Unknown dtype: {dtype}")
    if "H100" in properties.name:
        if dtype == torch.float32:
            return 67.5e12
        if dtype in (torch.bfloat16, torch.float16):
            return 1979e12 / 2
        raise ValueError(f"Unknown dtype: {dtype}")
    if "B200" in properties.name:
        if dtype == torch.float32:
            return 75e12
        if dtype in (torch.bfloat16, torch.float16):
            return 4.5e15 / 2
        raise ValueError(f"Unknown dtype: {dtype}")
    return None
\end{minted}

\subsubsection*{fp32}

Let's examine memory usage of these tensors. Memory is determined by the (i) number of values and (ii) data type of each value.

\begin{minted}{python}
    x = torch.zeros(4, 8)  
    assert x.dtype == torch.float32  # Default type
    assert x.numel() == 4 * 8
    assert x.element_size() == 4  # Float is 4 bytes
    assert get_memory_usage(x) == 4 * 8 * 4  # 128 bytes
\end{minted}

One matrix in the feedforward layer of GPT-3:

\begin{minted}{python}
     assert get_memory_usage(
        torch.empty(12288 * 4, 12288)) == 2304 * 1024 * 1024
\end{minted}

\subsubsection*{fp16}

The fp16 data type (also known as float16 or half precision) cuts down the memory.

\begin{minted}{python}
    x = torch.zeros(4, 8, dtype=torch.float16)  
    assert x.element_size() == 2
\end{minted}

However, the dynamic range (especially for small numbers) isn't great.

\begin{minted}{python}
    x = torch.tensor([1e-8], dtype=torch.float16)  
    assert x == 0  # Underflow!
\end{minted}

If this happens when you train, you can get instability.

\subsubsection*{bf16}

Google Brain developed brain floating point (bf16) in 2018 to address this issue. bf16 uses the same memory as fp16 but has the same dynamic range as fp32! The only catch is that the resolution is worse, but this matters less for deep learning.

\begin{minted}{python}
    x = torch.tensor([1e-8], dtype=torch.bfloat16)  
    assert x != 0  # No underflow!
\end{minted}

\subsubsection*{Mixed precision}

Implications on training:

\begin{itemize}
    \item Training with fp32 works, but requires lots of memory.
    \item Training with fp16 and even bf16 is risky, and you can get instability.
\end{itemize}

Solution: mixed precision training [Micikevicius+ 2017].

\begin{itemize}
    \item Use bf16 for parameters, activations, and gradients.
    \item Use fp32 for optimizer states.
\end{itemize}

Pytorch has an automatic mixed precision (AMP) library. Tries to cast things into bf16 when safe (matmuls, not exp).

\begin{minted}{python}
    with torch.amp.autocast("cuda", dtype=torch.bfloat16):
        x = torch.zeros(4, 8) 
\end{minted}

\subsubsection*{fp8}

In 2022, fp8 was standardized, motivated by machine learning workloads primer. H100s support two variants of FP8: E4M3 (range [-448, 448]) and E5M2 ([-57344, 57344]).

\subsubsection*{fp4}

In 2025, NVIDIA developed nvfp4. Only 4 bits per value! Values: -6, -4, -3, -2, -1.5, -1.0, -0.5, 0.0, 0.5, 1.0, 1.5, 2, 3, 4, 6. Use a separate scale factor per block, so actually get more dynamic range (but just can't vary freely from neighbors). Nemotron 3 Super was trained in NVFP4. Some of this is done in NVIDIA libraries outside of user control.

\subsection*{Tensors on GPUs}

By default, tensors are stored in CPU memory.

\begin{minted}{python}
    x = torch.zeros(32, 32)
    assert x.device == torch.device("cpu")
\end{minted}

However, what about GPUs?

\begin{minted}{python}
    device = cuda_if_available()
\end{minted}

In order to take advantage of the massive parallelism of GPUs, we need to move them to GPU memory.

\begin{minted}{python}
    x = x.to(device)
\end{minted}

Or create the tensor directly on the GPU:

\begin{minted}{python}
    with torch.device(device):
        x = torch.zeros(32, 32)
        assert x.device == device
\end{minted}

\subsection*{Tensor Einops}

Einops is a library for manipulating tensors where dimensions are named.
It is inspired by Einstein summation notation (Einstein, 1916).

\subsubsection*{Einops Motivation}

Traditional PyTorch code:

\begin{minted}{python}
    x = torch.ones(2, 2, 3)      # batch seq hidden  
    y = torch.ones(2, 2, 3)      # batch seq hidden  
    z = x @ y.transpose(-2, -1)  # batch seq seq 
\end{minted} 

Easy to mess up the dimensions (what is \(-2\), \(-1\)?)...

\subsubsection*{Einops Einsum}

Einsum is generalized matrix multiplication with good bookkeeping.

\begin{minted}{python}
    x = torch.ones(3, 4)  # seq1 hidden 
    y = torch.ones(4, 3)  # hidden seq2 
    
    # Old way
    z = x @ y   # seq1 seq2 
    
    # New (einops) way
    z = einsum(x, y, "seq1 hidden, hidden seq2 -> seq1 seq2")
\end{minted}

Let's try a more complex example...

\begin{minted}{python}
    x = torch.ones(2, 3, 4)  # batch seq1 hidden 
    y = torch.ones(2, 3, 4)  # batch seq2 hidden 
    
    # Old way
    z = x @ y.transpose(-2, -1)  # batch seq1 seq2  
    
    # New (einops) way
    z = einsum(x, y, 
        "batch seq1 hidden, batch seq2 hidden -> batch seq1 seq2") 
\end{minted}

Dimensions that are not named in the output are summed over.

\begin{minted}{python}
    # Or can use `...` to represent broadcasting 
    z = einsum(x, y, 
        "... seq1 hidden, ... seq2 hidden -> ... seq1 seq2")
\end{minted}

\subsubsection*{Einops Reduce}

You can reduce a single tensor via some operation (e.g., sum, mean, max, min).

\begin{minted}{python}
    x = torch.ones(2, 3, 4)  # batch seq hidden 
    
    # Old way
    y = x.sum(dim=-1)  
    
    # New (einops) way
    y = reduce(x, "... hidden -> ...", "sum")  
\end{minted}

\subsubsection*{Einops Rearrange}

Sometimes, a dimension represents two dimensions and you want to operate on one of them.

\begin{minted}{python}
    x = torch.ones(3, 8)  # seq total_hidden 
    w = torch.ones(4, 4)  # hidden1 hidden2 
    
    # Break up into two dimensions
    x = rearrange(x, 
        "... (heads hidden1) -> ... heads hidden1", heads=2)  
    # Perform the transformation by `w`
    x = einsum(x, w, 
        "... hidden1, hidden1 hidden2 -> ... hidden2")  
    
    # Combine `heads` and `hidden2` back together
    x = rearrange(x, 
        "... heads hidden2 -> ... (heads hidden2)")  
\end{minted}

\subsection*{Tensor FLOPs}

Having gone through all the operations, let us examine their computational cost. \\

A floating-point operation (FLOP) is a basic operation like addition (\(x + y\)) or multiplication (\(x \cdot y\)). \\

Two terribly confusing acronyms (pronounced the same!):

\begin{itemize}
    \item FLOPs: floating-point operations (measure of computation done).
    \item FLOP/s: floating-point operations per second (also written as FLOPS), which is used to measure the speed of hardware.
\end{itemize}

\subsubsection*{Intuitions}

\begin{itemize}
    \item Training GPT-3 (2020) took 3.14e23 FLOPs. 
    \item Training GPT-4 (2023) is speculated to take 2e25 FLOPs.
\end{itemize}

H100 has a peak performance of 1979 teraFLOP/s with sparsity, 50 without

\begin{minted}{python}
    h100_flop_per_sec = 1979e12 / 2
\end{minted}

8 H100s for 2 weeks:

\begin{minted}{python}
    total_flops = 8 * 2 * (60 * 60 * 24 * 7) * h100_flop_per_sec
\end{minted}

\subsubsection*{Linear model}

\begin{minted}{python}
    if torch.cuda.is_available():
        B = 16384  # Number of points
        D = 32768  # Dimension of each point
        K = 8192   # Number of outputs
    else:
        B = 1024
        D = 256
        K = 64
    x = torch.ones(B, D, device=cuda_if_available())
    w = torch.randn(D, K, device=cuda_if_available())
    y = x @ w
\end{minted}

How many FLOPs is this matmul? We have one multiplication \((x[i][j] \cdot w[j][k])\) and one addition per \((i, j, k)\) triple.

\begin{minted}{python}
    actual_num_flops = 2 * B * D * K  
\end{minted}

We can also time this operation to see how long it takes.

\begin{minted}{python}
    actual_time = benchmark(lambda: x @ w) 
\end{minted}

The actual FLOP/s of this operation:

\begin{minted}{python}
    actual_flop_per_sec = actual_num_flops / actual_time
\end{minted}

Each GPU has a specification sheet that provides the peak performance. Note that the FLOP/s depends heavily on the data type!

\begin{minted}{python}
    promised_flop_per_sec = get_promised_flop_per_sec(x.dtype)
\end{minted}

\subsubsection*{Model FLOPs utilization (MFU)}

\begin{definition}
    MFU = (actual FLOP/s) / (promised FLOP/s) [ignore communication/overhead].
\end{definition}

\begin{minted}{python}
     mfu = actual_flop_per_sec / promised_flop_per_sec 
\end{minted}

Usually, MFU of \(\geq\) 0.5 is quite good! But why is MFU not closer to 1? To answer this question, we need to look more closely at how computations are done on GPUs...

\subsection*{Arithmetic Intensity}

How to compute a thing:

\begin{enumerate}
    \item Send inputs from memory to accelerator.
    \item Perform computation.
    \item Send outputs from accelerator to memory.
\end{enumerate}

How long does this take? Depends on two things:

\begin{enumerate}
    \item Accelerator speed (FLOP/s)
    \item Memory bandwidth (bytes/s)
\end{enumerate}

\begin{minted}{python}
    assert h100_flop_per_sec == 1979e12 / 2
    assert h100_bytes_per_sec == 3.35e12
\end{minted}

\subsubsection*{RELUs}

\begin{minted}{python}
    n = 1024 * 1024
    x = torch.ones(n, dtype=torch.bfloat16, device=cuda())
    y = torch.relu(x)
    
    bytes = (2 * n) + (2 * n)  # Read x, write y
    flops = n  # n comparisons
    
    communication_time = bytes / h100_bytes_per_sec  
    computation_time = flops / h100_flop_per_sec  
\end{minted}

Assume we can overlap communication and computation perfectly.

\begin{minted}{python}
    total_time = max(communication_time, computation_time)  
\end{minted}

What is the bottleneck?

\begin{itemize}
    \item Memory-bound: communication time \(>\) computation time.
    \item Compute-bound: computation time \(>\) communication time.
\end{itemize}

In this case, ReLU is memory-bound. \\

Accelerator intensity: how much work can the accelerator do per byte transferred?

\begin{minted}{python}
    accelerator_intensity = h100_flop_per_sec / h100_bytes_per_sec  
\end{minted}

Arithmetic intensity: how much actual work per byte for this workload?

\begin{minted}{python}
    arithmetic_intensity = flops / bytes  # 1/4
\end{minted}

What is the bottleneck?

\begin{itemize}
    \item Memory-bound: arithmetic intensity \(<\) accelerator intensity.
    \item Compute-bound: arithmetic intensity \(>\) accelerator intensity.
\end{itemize}

\begin{minted}{python}
    assert arithmetic_intensity < accelerator_intensity
\end{minted}

In general, we'll find ourselves memory bound. Can we increase arithmetic intensity?

\subsubsection*{GELUs}

\begin{minted}{python}
    n = 1024 * 1024
    x = torch.ones(n, dtype=torch.bfloat16, device=cuda())
    y = F.gelu(x) 
    
    bytes = (2 * n) + (2 * n)  # Read x, write y
    flops = 20 * n  # tanh can be approximated
    
    arithmetic_intensity = flops / bytes  
    accelerator_intensity = h100_flop_per_sec / h100_bytes_per_sec 
    
    assert arithmetic_intensity < accelerator_intensity
\end{minted}

Note that GeLU does more work than ReLU per byte moved, so it has higher arithmetic intensity. But still memory-bound! \\

In other words, ReLU is not faster than GeLU (when doing things in an isolated way).

\subsubsection*{Dot Products}

\begin{minted}{python}
    n = 1024 * 1024
    x = torch.ones(n, dtype=torch.bfloat16, device=cuda())
    w = torch.ones(n, dtype=torch.bfloat16, device=cuda())
    y = x @ w
    
    bytes = (2 * n) + (2 * n) + 2  # Read x, read w, write y
    flops = 2 * n - 1  # n multiplications, n-1 additions
    
    arithmetic_intensity = flops / bytes  # 1/2 
    accelerator_intensity = h100_flop_per_sec / h100_bytes_per_sec  
    
    assert arithmetic_intensity <accelerator_intensity
\end{minted}

Memory-bound!

\subsubsection*{Matrix-Vector Products}

\begin{minted}{python}
    n = 1024
    x = torch.ones(n, dtype=torch.bfloat16, device=cuda())
    w = torch.ones(n, n, dtype=torch.bfloat16, device=cuda())
    y = x @ w
    
    bytes = (2 * n) + (2 * n * n) + (2 * n)
    flops = n * (2 * n - 1)  # n dot-products
    
    arithmetic_intensity = flops / bytes  # 1 
    accelerator_intensity = h100_flop_per_sec / h100_bytes_per_sec 
    
    assert arithmetic_intensity < accelerator_intensity
\end{minted}

Memory-bound!

\subsubsection*{MatMuls}

\begin{minted}{python}
    n = 1024
    x = torch.ones(n, n, dtype=torch.bfloat16, device=cuda())
    w = torch.ones(n, n, dtype=torch.bfloat16, device=cuda())
    y = x @ w
    
    bytes = (2 * n * n) + (2 * n * n) + (2 * n * n)
    flops = n * n * (2 * n - 1)  # n^2 dot products
    
    arithmetic_intensity = flops / bytes  # n/3 
    accelerator_intensity = h100_flop_per_sec / h100_bytes_per_sec  
    
    assert arithmetic_intensity > accelerator_intensity
\end{minted}

Finally, compute-bound! 

\begin{itemize}
    \item As long as we have large matrices, we're compute-bound (saturating the accelerator).
    \item Training Transformers involves big matrix multiplications.
    \item Matrix-vector product is what happens during inference, which is why inference is memory-bound.
\end{itemize}

Note: arithmetic/accelerator intensity also depends on the precision (bf16 versus fp32).

\subsection*{Roofline Plots}

We can visualize the relationship between arithmetic intensity and performance using roofline plots.

\begin{itemize}
    \item Each slice on the x-axis is a particular computation (with some arithmetic intensity).
    \item Each piecewise linear function corresponds to a particular hardware.
    \item Kink is the accelerator intensity (transition from memory-bound to compute-bound).
\end{itemize}

We can now relate this back: MFU = \(\min\)(1, arithmetic-intensity / accelerator-intensity)

\subsection*{Gradient Basics}

So far, we've constructed tensors and passed them through operations (forward). Now, we're going to compute the gradient (backward). As a simple example, let's consider the simple linear model:

\[
    y = 0.5 (x \cdot w - 5)^2
\]

\textbf{Forward pass:} compute loss

\begin{minted}{python}
    x = torch.tensor([1., 2, 3])
    w = torch.tensor([1., 1, 1], requires_grad=True)
    pred_y = x @ w
    loss = 0.5 * (pred_y - 5).pow(2)
\end{minted}

\textbf{Backward pass:} compute gradients

\begin{minted}{python}
    loss.backward()
    assert torch.equal(w.grad, torch.tensor([1, 2, 3]))
\end{minted}

\subsection*{Gradient FLOPs}

Let us count the FLOPs for computing gradients.

\begin{minted}{python}
    B = 1024  # Number of points
    D = 256   # Dimension
\end{minted}

Define a simplified model (2-layer linear network):

\begin{minted}{python}
    x = torch.ones(B, D, device=cuda())
    w1 = torch.randn(D, D, device=cuda(), requires_grad=True)
    w2 = torch.randn(D, D, device=cuda(), requires_grad=True)
    
    # Forward pass
    h1 = einsum(x, w1, "batch in, in out -> batch out")
    h2 = einsum(h1, w2, "batch in, in out -> batch out")
    loss = (h2.mean() - 0)**2  # Regress everything to 0
    
    # Backward pass
    h1.retain_grad()  # For debugging
    h2.retain_grad()  # For debugging
    loss.backward()
\end{minted}

\subsubsection*{Zoom in on one layer}

Let's focus on the second layer (h2 = h1 @ w2). \\

\textbf{Forward pass:} Recall the number of forward FLOPs:  

\begin{minted}{python}
    num_forward_flops = 2 * B * D * D
\end{minted}

\textbf{Backward pass:} How many FLOPs is running the backward pass?

\begin{minted}{python}
    h1_grad = einsum(h2.grad, w2, "batch out, in out -> batch in")
    assert torch.allclose(h1.grad, h1_grad)
    
    w2_grad = einsum(h2.grad, h1, "batch out, batch in -> in out")
    assert torch.allclose(w2.grad, w2_grad)
    
    num_backward_flops = (2 * B * D * D) + (2 * B * D * D)
\end{minted}

Note that the backward pass is 2x more expensive than the forward pass.

\subsubsection*{Consider all layers}

This was just for w2, need to apply it to all parameters in the network. \\

Putting it together:

\begin{itemize}
    \item Forward pass: 2 (\# data points) (\# parameters) FLOPs.
    \item Backward pass: 4 (\# data points) (\# parameters) FLOPs.
    \item Total: 6 (\# data points) (\# parameters) FLOPs.
\end{itemize}

This is for multilayer perceptrons (MLPs) ... but it turns out to be a good approximation for Transformers for short context lengths as well.

\begin{minted}{python}
class Block(torch.Module):
    """Simple block that applies a linear transformation 
    followed by a ReLU nonlinearity."""
    def __init__(self, dim: int):
        super().__init__()
        self.weight = nn.Parameter(
            torch.randn(dim, dim) / math.sqrt(dim))
    
    def forward(self, x: torch.Tensor) -> torch.Tensor:
        x = x @ self.weight  # Linear
        x = F.relu(x)        # Activation
        return x

class DeepNetwork(torch.Module):
    """Map `dim`-vector to a `dim`-vector."""
    def __init__(self, dim: int, num_layers: int):
        super().__init__()
        self.layers = nn.ModuleList([
            Block(dim) for i in range(num_layers)])
    
    def forward(self, x: torch.Tensor) -> torch.Tensor:
        # Apply all the layers sequentially
        for layer in self.layers:
            x = layer(x)  
        return x

class AdaGrad(torch.optim.Optimizer):
    def __init__(self,  params: Iterable[nn.Parameter], lr: float):
        super(AdaGrad, self).__init__(params, dict(lr=lr))
    
    def step(self):
        for group in self.param_groups:
            lr = group["lr"]
            for p in group["params"]:
                state = self.state[p]
                grad = p.grad.data
                # Get squared gradients g2 = sum_{i<t} g_i^2
                g2 = state.get("g2", torch.zeros_like(grad))
                # Update optimizer state
                g2 += torch.square(grad)
                state["g2"] = g2
                # Update parameters
                p.data -= lr * grad / torch.sqrt(g2 + 1e-5)

class DeepNetworkCheckpointed(torch.Module):
    """Same as DeepNetwork, but with activation checkpointing."""
    def __init__(self, dim: int, num_layers: int):
        super().__init__()
        self.layers = nn.ModuleList([
            Block(dim) for i in range(num_layers)])
    
    def forward(self, x: torch.Tensor) -> torch.Tensor:
        # Apply all the layers sequentially
        for layer in self.layers:
            # KEY: only store activations at checkpoints
            x = torch.utils.checkpoint.checkpoint(layer, x)  
        return x
\end{minted}

\subsection*{Deep Network}

Consider a deep network with \(L\) layers and \(D\)-dimensional inputs, activations, and outputs.

\begin{minted}{python}
    # Define the network
    L = 3  # Number of layers
    D = 8  # Dimensionality of input, activations and output
    model = DeepNetwork(dim=D, num_layers=L).to(cuda())
    num_parameters = get_num_parameters(model)  
    assert num_parameters == (D * D) * L
    
    # Run the model on a batch of data
    B = 4  # Batch size
    x = torch.randn(B, D, device=cuda())  
    y = model(x)  
\end{minted}

\subsection*{Optimizer}

Let's define the AdaGrad optimizer:

\begin{itemize}
    \item momentum \(=\) SGD \(+\) exponential averaging of grad.
    \item AdaGrad \(=\) SGD \(+\) averaging by grad squared.
    \item RMSProp \(=\) AdaGrad but with exponential averaging of grad squared.
    \item Adam \(=\) RMSProp \(+\) momentum.
\end{itemize}

\begin{minted}{python}
    optimizer = AdaGrad(model.parameters(), lr=0.01)  
    state = model.state_dict()
    
    # Compute gradients
    x = torch.randn(B, D, device=cuda())
    y = torch.tensor([4., 5.], device=cuda())
    pred_y = model(x).mean()  
    loss = F.mse_loss(input=pred_y, target=y)
    loss.backward()
    
    # Take a step
    optimizer.step()

    # Free up the memory
    optimizer.zero_grad(set_to_none=True)
\end{minted}

\subsection*{Memory}

\begin{minted}{python}
    num_parameters = D * D * L
    parameter_memory = 2 * num_parameters  # (2 bytes for bf16) 
    gradient_memory = 2 * num_parameters  # (2 bytes for bf16) 
    optimizer_memory = 4 * num_parameters  # (4 bytes for fp32) 
    activation_memory = 2 * (B * D * L)  # (2 bytes for bf16) 
\end{minted}

It is customary to use fp32 for stability (accumulating averages over powers over many steps).

\begin{itemize}
    \item AdaGrad: 4 bytes/parameter for storing second moments.
    \item Adam: 8 bytes/parameter for storing first and second moments.
\end{itemize}

\subsection*{Compute (for one training step)}

\begin{minted}{python}
    num_parameters = D * D * L
    flops = 6 * B * num_parameters
\end{minted}

\subsection*{Transformer}

The accounting for a Transformer is more complicated, but the same idea.

\begin{minted}{python}
def train_loop():
    # True linear function with weights (0, 1, 2, ..., D-1)
    D = 16  # Dimensionality
    true_w = torch.arange(D, dtype=torch.float32, device=cuda())
    # Data loader that generates (x, y) pairs
    B = 4  # Batch size
    def get_batch() -> tuple[torch.Tensor, torch.Tensor]:
        x = torch.randn(B, D).to(cuda())
        true_y = x @ true_w
        return (x, true_y)
    # Define the model and optimizer
    L = 2  # Number of layers
    model = DeepNetwork(dim=D, num_layers=L).to(cuda()) 
    optimizer = AdaGrad(model.parameters(), lr=0.01) 
    # Train!
    num_train_steps = 10
    for t in range(num_train_steps):
        x, y = get_batch()
        # Forward (compute loss)
        pred_y = model(x).mean()  
        loss = F.mse_loss(pred_y, y)
        # Backward (compute gradients)
        loss.backward()
        optimizer.step()
        optimizer.zero_grad(set_to_none=True)
\end{minted}

\subsection*{Gradient Accumulation}

Large batch sizes: improve training stability. However, activation memory scales with batch size, so might run out.

\begin{minted}{python}
    B = 64     # Batch size
    D = 1024   # Dimensionality
    L = 16     # Number of layers
    activation_memory = 2 * B * D * L  # (2 bytes for bf16)
\end{minted}

Accumulation:

\begin{itemize}
    \item Compute gradient on micro batches.
    \item Accumulate the gradients (don't zero it out).
    \item Every batch\_size / micro\_batch\_size steps, update the parameters and zero out the gradients
\end{itemize}

\begin{minted}{python}
    micro_batch_size = 256
    activation_memory = 2 * micro_batch_size * D * L
\end{minted}

\subsection*{Activation Checkpointing}

For training, we need to store the activations of all layers. For inference, we do not compute gradients, so we only need to store the current layer's activations. The memory usage is

\begin{minted}{python}
    B = 64     # Batch size
    D = 1024   # Dimensionality
    L = 16     # Number of layers
    
    x = torch.randn(B, D, device=cuda(), requires_grad=True)
    activation_memory = 2 * B * D * L  
    
    model = DeepNetwork(dim=D, num_layers=L).to(cuda())  
    memory = get_memory_usage(lambda: model(x).sum().backward()) 
\end{minted}

Can we reduce this? Activation checkpointing \(=\) gradient checkpointing \(=\) rematerialization.

\begin{itemize}
    \item Forward pass: keep only activations at subset of layers.
    \item Backward pass: recompute the missing activations from the last checkpoint.
\end{itemize}

Philosophy: tradeoff memory for compute.

\begin{minted}{python}
    # Store all activations:    x g1 h1 g2 h2 g3 h3 g4 h4
    # Activation checkpointing: x    h1    h2    h3    h4
    
    # Define the model with checkpointing
    model = DeepNetworkCheckpointed(dim=D, num_layers=L).to(cuda())  
    checkpointed_memory = get_memory_usage()
\end{minted}

Can we reduce this even more, especially for deep networks (large \(L\))?

\begin{minted}{python}
    # Store all layers:   | h1 h2 h3 h4 h5 h6 h7 h8 h9 |
    # Store no layers:    |                            |
    # Store some layers:  |    h3       h6          h9 |
\end{minted}

How frequently to checkpoint?

\begin{itemize}
    \item If store each layer's activations, then activation memory is \(O(L)\) and no recomputation.
    \item If store no activations, then activation memory is \(O(1)\) and compute is \(O(L^2)\) (recompute from the start for each layer).
    \item If store every \(\sqrt{L}\) layers, then activation memory is \(O(\sqrt{L})\) and \(O(L)\) recomputation.
\end{itemize}

\begin{minted}{python}
def get_num_parameters(model: torch.Module) -> int:
    return sum(param.numel() for param in model.parameters())
\end{minted}

\end{document}
